% !TEX root = ../main.tex
\chapter{ケーススタディ：lossy migration と仕様弱化}
\label{chap:ch23_lossy_migration}

\begin{chapterabstract}
本章では、情報を失うマイグレーションを扱う。前章の lossless migration では、旧形式と新形式の間に両方向の round-trip law を置くことができた。しかし実務では、すべての変更が同型になるわけではない。代表例は、旧スキーマの \texttt{fullName} を、新スキーマの \texttt{firstName} と \texttt{lastName} に分ける変更である。空白、ミドルネーム、敬称、文化圏ごとの姓名順、表記ゆれがあるため、完全復元は一般には成立しない。

このような変更で重要なのは、「戻せないから検証できない」と諦めることではない。完全復元という強い仕様を捨て、証明できる弱い仕様へ落とす。たとえば、正規化後の表示が一致する、well-formed な入力に限れば戻せる、新システムで必要な公開表示だけは保存される、といった主張である。本章では、この仕様弱化を Lean 4 の小さなモデルで表す。
\end{chapterabstract}

\begin{learninggoals}
\item 情報損失を伴う移行が完全同型ではない理由を説明できる。
\item 完全な round-trip、条件付き round-trip、正規化後一致、表示互換を区別できる。
\item well-formed 条件を Lean の述語として導入できる。
\item 証明できない仕様を、証明できる弱い仕様へ落とす手順を理解する。
\item lossy migration の保証範囲を設計文書に明示できる。
\end{learninggoals}

\section{この章を読む理由}

マイグレーションのレビューでは、「互換性があります」という言葉がよく使われる。しかし、その意味は曖昧である。旧データを新形式で読めるという意味なのか。新形式から旧形式へ戻せるという意味なのか。画面表示だけが同じという意味なのか。検索キーや集計結果だけが同じという意味なのか。これらは同じではない。

\begin{intuitionbox}
lossy migration では、元の情報の一部が新形式に入らない。したがって、「完全に戻る」という仕様は強すぎる。必要なのは、業務上守りたい観測を選び、その観測については保存される、と言い直すことである。
\end{intuitionbox}

本章の例では、旧形式が \texttt{fullName} を持ち、新形式が \texttt{firstName} と \texttt{lastName} だけを持つと考える。実務では \texttt{fullName} は単なる文字列であることが多い。しかし文字列パーサの詳細に入ると本章の主題がぼやけるため、Lean では「トークン化済みの \texttt{fullName}」として、\texttt{first}、任意の \texttt{middle}、\texttt{last} を持つ小さなモデルを使う。これは、生の文字列に含まれていた情報の一部が新形式に収まらない状況を表すためのモデルである。

\FigurePlaceholder{fig-ch23-lossy-overview.png}{0.92}{lossy migration では完全復元ではなく弱めた仕様を証明する}{fig:ch23-lossy-overview}

\section{情報が欠ける移行は同型ではない}

第22章では、旧形式と新形式が同じ情報を持つ場合を扱った。そのときは、\texttt{migrate : Old \textrightarrow New} と \texttt{rollback : New \textrightarrow Old} を定義し、次の二つを証明できた。

\begin{itemize}
\item \texttt{rollback (migrate old) = old}
\item \texttt{migrate (rollback new) = new}
\end{itemize}

これは、旧形式と新形式が同型であることを表していた。一方、lossy migration では少なくとも片方の等式が壊れる。

\begin{definitionbox}
本章では、\textbf{lossy migration} を「旧形式から新形式へ移すとき、旧形式の情報の一部が新形式に保存されない移行」と呼ぶ。lossy migration は完全な同型ではない。したがって、完全な round-trip law を無条件に要求してはいけない。
\end{definitionbox}

たとえば、旧形式がミドルネームを含み、新形式が \texttt{firstName} と \texttt{lastName} しか持たない場合、\texttt{middle} は保存先を持たない。新形式から旧形式へ戻すとき、\texttt{middle = none} と埋めることはできる。しかし、もともと \texttt{some "Byron"} だったか、\texttt{some "Augusta"} だったか、\texttt{none} だったかは復元できない。

\begin{warningbox}
\texttt{migrate} と \texttt{rollback} の二つの関数を書けたとしても、それだけでは安全な相互変換ではない。安全な相互変換だと言うには、戻した結果が本当に元と等しいことを証明する必要がある。
\end{warningbox}

\section{例：\texttt{fullName} を \texttt{firstName} / \texttt{lastName} に分ける}

まず、Lean で小さなモデルを作る。旧スキーマ \texttt{UserV1} は \texttt{fullName} を持つ。新スキーマ \texttt{UserV2} は \texttt{firstName} と \texttt{lastName} だけを持つ。\texttt{FullName} の \texttt{middle : Option String} は、旧形式には存在したが新形式には入りきらない情報を表している。

\begin{leanbox}
ここでは実際の文字列分割関数ではなく、分割済みの名前モデルを使う。これにより、文字列処理の細部ではなく、「新形式に入らない情報があると round-trip が壊れる」という構造だけを検証できる。
\end{leanbox}

\begin{listing}[htbp]
\begin{minted}{lean4}
namespace Ch23

structure FullName where
  first : String
  middle : Option String
  last : String
deriving Repr, DecidableEq

structure UserV1 where
  id : Nat
  fullName : FullName
deriving Repr, DecidableEq

structure UserV2 where
  id : Nat
  firstName : String
  lastName : String
deriving Repr, DecidableEq

def migrate (u : UserV1) : UserV2 :=
  { id := u.id,
    firstName := u.fullName.first,
    lastName := u.fullName.last }

def rollback (v : UserV2) : UserV1 :=
  { id := v.id,
    fullName := { first := v.firstName,
                  middle := none,
                  last := v.lastName } }
\end{minted}
\caption{旧形式・新形式・移行関数}
\label{lst:ch23-model}
\end{listing}

この \texttt{migrate} は、旧形式の \texttt{middle} を新形式へ保存しない。したがって、旧データから新データへ移し、再び旧データへ戻す方向は一般には元に戻らない。

ただし、逆方向は戻る。新形式には最初から \texttt{middle} がないので、\texttt{rollback} で \texttt{middle = none} を補い、もう一度 \texttt{migrate} すると、元の新形式に戻る。

\begin{listing}[htbp]
\begin{minted}{lean4}
theorem new_round_trip (v : UserV2) :
    migrate (rollback v) = v := by
  cases v with
  | mk id firstName lastName =>
      rfl
\end{minted}
\caption{新形式については round-trip が成立する}
\label{lst:ch23-new-roundtrip}
\end{listing}

ここで得られた性質は「新形式に入る情報だけを見れば安定している」という主張である。これは有用だが、旧形式の情報をすべて保存しているという意味ではない。

\section{完全な round-trip が壊れることを明示する}

lossy migration では、壊れる仕様を曖昧にしたままにしないことが重要である。Lean では、具体例を使って「強い仕様は成立しない」ことも明示できる。

\begin{listing}[htbp]
\begin{minted}{lean4}
def oldWithMiddle : UserV1 :=
  { id := 42,
    fullName :=
      { first := "Ada",
        middle := some "Byron",
        last := "Lovelace" } }

#eval migrate oldWithMiddle  -- => { id := 42, firstName := "Ada", lastName := "Lovelace" }
#eval rollback (migrate oldWithMiddle)  -- => { id := 42, fullName := { first := "Ada", middle := none, last := "Lovelace" } }

theorem oldWithMiddle_not_restored :
    rollback (migrate oldWithMiddle) ≠ oldWithMiddle := by
  decide

def StrongRoundTrip : Prop :=
  ∀ u : UserV1, rollback (migrate u) = u

theorem strong_spec_fails : ¬ StrongRoundTrip := by
  intro h
  exact oldWithMiddle_not_restored (h oldWithMiddle)
\end{minted}
\caption{ミドルネームを含む旧データでは完全復元できない}
\label{lst:ch23-strong-fails}
\end{listing}

\texttt{StrongRoundTrip} は、前章の lossless migration なら目標にできた仕様である。しかし本章の移行では、\texttt{oldWithMiddle} が反例になる。証明の意味は単純である。もし全旧データについて完全復元できるなら、\texttt{oldWithMiddle} も復元できるはずである。しかし Lean は、この具体例が復元されないことを計算で確認できる。したがって、強い仕様は成り立たない。

\begin{softwarebox}
この段階でレビューに書くべきことは、「この移行は不可逆である」という事実である。不可逆であること自体が失敗なのではない。不可逆なのに可逆であるかのように扱うことが失敗である。
\end{softwarebox}

\section{仕様弱化：証明可能な性質に落とす}

完全復元が成り立たないなら、次に行うべきことは、業務上本当に必要な性質を選ぶことである。たとえば、管理画面ではミドルネームを表示しない方針かもしれない。検索や通知では、\texttt{firstName} と \texttt{lastName} だけが使われるかもしれない。その場合、保存すべき性質は「旧データ全体が復元できること」ではなく、「公開表示が一致すること」である。

\begin{definitionbox}
\textbf{仕様弱化}とは、証明できない強い仕様を、業務上必要で証明可能な弱い仕様へ置き換えることである。本章の例では、\texttt{rollback (migrate u) = u} を諦め、\texttt{displayV2 (migrate u) = canonicalDisplayV1 u} を証明する。
\end{definitionbox}

\begin{listing}[htbp]
\begin{minted}{lean4}
def canonicalDisplayV1 (u : UserV1) : String :=
  u.fullName.first ++ " " ++ u.fullName.last

def displayV2 (v : UserV2) : String :=
  v.firstName ++ " " ++ v.lastName

def WeakDisplaySpec : Prop :=
  ∀ u : UserV1, displayV2 (migrate u) = canonicalDisplayV1 u

theorem weak_display_spec_holds : WeakDisplaySpec := by
  intro u
  rfl
\end{minted}
\caption{完全復元ではなく表示互換を証明する}
\label{lst:ch23-weak-display}
\end{listing}

この定理は、すべての旧データについて、新形式に移した後の表示が、旧形式を公開表示用に正規化した結果と一致することを述べている。ミドルネームを含む旧データでも成立する。ただし、それはミドルネームが保存されるという意味ではない。ミドルネームを表示仕様の外に置く、という仕様判断をしたから成立する。

\FigurePlaceholder{fig-ch23-spec-weakening.png}{0.92}{強い仕様を弱めて証明可能な仕様へ落とす流れ}{fig:ch23-spec-weakening}

仕様弱化は、単なる妥協ではない。保証範囲を明確にする作業である。強い仕様が成り立たないときに、どの観測なら保存できるのかを明示することで、レビューは具体的になる。

\begin{center}
\begin{tabular}{p{0.25\linewidth}p{0.34\linewidth}p{0.31\linewidth}}
\toprule
仕様の種類 & 形 & 本章の例での状態 \\
\midrule
完全復元 & \texttt{rollback (migrate u) = u} & ミドルネームがあると壊れる \\
条件付き復元 & \texttt{NoMiddle u \textrightarrow rollback (migrate u) = u} & ミドルネームがない旧データなら成立 \\
正規化後一致 & \texttt{normalizeV1 (normalizeV1 u) = normalizeV1 u} & 一度丸めた後は安定する \\
表示互換 & \texttt{displayV2 (migrate u) = canonicalDisplayV1 u} & 全旧データで成立 \\
\bottomrule
\end{tabular}
\end{center}

\section{well-formed 条件を導入する}

仕様弱化には二つの方向がある。一つは、保存する観測を弱めることである。前節の表示互換がこれにあたる。もう一つは、入力範囲を狭めることである。旧データのうち、ミドルネームを持たないものだけを対象にすれば、完全復元は成立する。

\begin{definitionbox}
\textbf{well-formed 条件}とは、証明したい性質が成立する入力の範囲を明示する述語である。本章のモデルでは、\texttt{NoMiddle u} を「旧データ \texttt{u} はミドルネームを持たない」という条件として使う。
\end{definitionbox}

\begin{listing}[htbp]
\begin{minted}{lean4}
def NoMiddle (u : UserV1) : Prop :=
  u.fullName.middle = none

theorem rollback_after_migrate_if_no_middle
    (u : UserV1) (h : NoMiddle u) :
    rollback (migrate u) = u := by
  cases u with
  | mk id fullName =>
      cases fullName with
      | mk first middle last =>
          simp [NoMiddle] at h
          cases h
          rfl
\end{minted}
\caption{well-formed な旧データに限れば完全復元できる}
\label{lst:ch23-wellformed}
\end{listing}

この証明では、\texttt{NoMiddle u} という前提を使う。前提があるため、旧データの \texttt{middle} は \texttt{none} だとわかる。すると、\texttt{rollback} が補う \texttt{middle = none} と一致し、構造体全体が元に戻る。

\begin{softwarebox}
実務では、well-formed 条件は「このバッチ対象にはミドルネームが存在しない」「この国コードのユーザーだけを移行する」「過去に正規化済みのデータだけを対象にする」のような前提になる。Lean で定理に前提として現れる条件は、運用上も確認すべき条件である。
\end{softwarebox}

\section{正規化として見る}

lossy migration は、旧データを新形式で表せる範囲へ丸める操作として見ることもできる。本章の例では、\texttt{rollback (migrate u)} が「ミドルネームを落とした旧データ」を作る。これを正規化と呼ぶ。

\begin{intuitionbox}
正規化は、「情報を完全に戻す」操作ではない。むしろ、「このシステムが今後扱う標準形へそろえる」操作である。一度標準形へそろえた後は、同じ正規化を繰り返しても結果は変わらない。
\end{intuitionbox}

\begin{listing}[htbp]
\begin{minted}{lean4}
def normalizeV1 (u : UserV1) : UserV1 :=
  rollback (migrate u)

theorem migrate_after_normalize (u : UserV1) :
    migrate (normalizeV1 u) = migrate u := by
  cases u with
  | mk id fullName =>
      cases fullName with
      | mk first middle last =>
          rfl

theorem normalize_idempotent (u : UserV1) :
    normalizeV1 (normalizeV1 u) = normalizeV1 u := by
  cases u with
  | mk id fullName =>
      cases fullName with
      | mk first middle last =>
          rfl
\end{minted}
\caption{lossy migration を正規化として見る}
\label{lst:ch23-normalize}
\end{listing}

\texttt{migrate\_after\_normalize} は、正規化してから移行しても、最初に移行した結果と変わらないことを述べる。\texttt{normalize\_idempotent} は、正規化を二度かけても一度かけた結果と同じであることを述べる。これは lossy migration を扱うときの典型的な保証である。

圏論的には、ここには商や抽象化の発想がある。複数の旧データが、同じ新データへ送られる。\texttt{middle} だけが違う旧データは、新形式では区別されない。同じ新データへ落ちる旧データたちを「同じものとして扱う」と考えると、lossy migration は旧世界からより粗い世界への写像として読める。

\FigurePlaceholder{fig-ch23-normalization-quotient.png}{0.92}{複数の旧データが同じ新データへ落ちる様子}{fig:ch23-normalization-quotient}

\section{仕様レビューでの使い方}

lossy migration のレビューでは、コードだけでなく仕様の強さを確認する。次の問いを順に書き出すとよい。

\begin{softwarebox}
\begin{itemize}
\item 旧形式と新形式の型は何か。
\item 新形式に保存されない情報は何か。
\item 完全な round-trip は成り立つか。成り立たないなら反例は何か。
\item well-formed 条件を置けば完全復元できるか。
\item 入力範囲を狭めない場合、どの観測なら保存されるか。
\item 正規化後の値は安定しているか。
\item 証明した性質と、証明していない性質を設計文書に分けて書いているか。
\end{itemize}
\end{softwarebox}

この章の例でいえば、設計文書には次のように書ける。

\begin{quote}
本移行は、\texttt{middle} に相当する情報を新スキーマに保存しないため、旧データ全体に対する完全復元は保証しない。ただし、\texttt{middle = none} の旧データについては \texttt{rollback (migrate u) = u} が成立する。また、全旧データについて、\texttt{first + last} による公開表示は移行後も保存される。移行後に旧形式へ戻す場合、値は \texttt{middle = none} の正規化済み旧形式になる。
\end{quote}

このように書くと、「互換性あり」という曖昧な言葉ではなく、何が保存され、何が保存されないかをレビューできる。

\section{本章ではここまで}

本章では、\texttt{fullName} から \texttt{firstName} / \texttt{lastName} への移行を題材に、lossy migration と仕様弱化を扱った。実務の文字列処理そのものは、もっと多くの例外を持つ。姓名順、敬称、法人名、ニックネーム、複数空白、Unicode 正規化なども考える必要がある。本章の Lean モデルは、それらを直接検証しているわけではない。

それでも、この小さなモデルには価値がある。どの仕様が強すぎるのか、どの条件なら復元できるのか、どの観測なら保存できるのかを分けて書けるからである。

\begin{warningbox}
本章の証明は、実装の文字列パーサが正しいことを保証しない。保証しているのは、モデル化した範囲での情報損失、条件付き復元、表示互換、正規化安定性である。実システムに適用する場合は、パーサ、Unicode 正規化、入力データ品質、運用上の対象範囲を別途確認する必要がある。
\end{warningbox}

\section{本章のまとめ}

lossy migration は完全同型ではない。二つの変換関数があっても、round-trip law がなければ安全な相互変換とは言えない。

強い仕様が成り立たないときは、反例を明示する。これにより、保証できないことを設計レビューで隠さず扱える。

仕様弱化では、完全復元ではなく、表示互換、検索キー保存、集計結果保存など、業務上必要な観測を選んで証明する。

well-formed 条件を置けば、入力範囲を狭めた条件付き round-trip を証明できる場合がある。

lossy migration は正規化として見ることもできる。一度正規化した値が安定することは、実務上有用な保証になる。

% The namespace is closed in the extracted Lean file.
